\lecture{2}{Sat 27 Aug 2022 13:14}{Completeness of $\mathbb{R}$ }
\subsubsection{Archimedean Property}

\begin{theorem}[Archimedean property]
	For every $x \in \mathbb{R}$ there is $n \in  \mathbb{N}$ such that $x<n$.
\end{theorem}
\begin{proof}
Suppose for the sake of contradiction that there exists some $x \in \mathbb{R}$ that bound $\mathbb{N}$ from above, i.e. $n\le x$ for any $n \in  \mathbb{N}$.

Let $u = \operatorname{sup} \mathbb{N}$. Hence there exists $n_1 \in \mathbb{N}$ such that $u -1 < n_1\le u$. But $n+1 \in \mathbb{N}$ and $n+1 > u$, so  $u$ is not an upper bound of  $\mathbb{N}$.
\end{proof}

\begin{corollary}
	For any $y>0, x \in \mathbb{R}$ there exists $n \in \mathbb{N}$ such that $ny>x$.
\end{corollary}
\begin{proof}
There exists $n > \frac{x}{y}$.
\end{proof}

\begin{corollary}
	For any $x>0$ there is $n \in \mathbb{N}$ such that $0<\frac{1}{n}<x$.
\end{corollary}
\begin{proof}
There exists $n>\frac{1}{x}$.
\end{proof}

\begin{corollary}[of the previous corollary]
	$\operatorname{inf} \{\frac{1}{n}:n \in \mathbb{N}\} =0$.
\end{corollary}
\begin{proof}
Verify the definition of infimum.
\begin{enumerate}
	\item $0\le \frac{1}{n}$ for $n \in \mathbb{N}$.
	\item For any $v>0$, by the previous corollary,
		exists  $n$ such that  $0<\frac{1}{n}<v$.
\end{enumerate}
\end{proof}


\begin{theorem}[Density of rationals in the reals.]
	For any $a , b \in \mathbb{R}$, $a<b$, there exists $r \in \mathbb{Q}$ such that  $a<r<b$.
\end{theorem}
\begin{proof}
Assume first that $a$ is a positive number. Choose step size $y = \frac{1}{n} < b-a$. By the first corollary there exists $m \in \mathbb{N}$ such that  $a< \frac{m}{n}$.

Let $m$ be smallest such $m \in \mathbb{N}$ (well-ordered property of $\mathbb{N}$). Then $\frac{m -1}{n}\le a < \frac{m}{n}$. Now $\frac{m}{n}\le a+\frac{1}{n}<b$.

If $a<0$ and  $b>0$ then  $0 = r \in \mathbb{Q}$. If $a<0$ and  $b\le 0$ then we flip sign. There exists $r \in \mathbb{Q}$ such that $-b <r < -a$ and hence  $a<-r<b$. 

\begin{figure}[ht]
    \centering
    \incfig{density-of-rationals}
    \caption{density-of-rationals}
    \label{fig:density-of-rationals}
\end{figure}
\end{proof}


\begin{theorem}[Existence of $\sqrt{2}$]
	There exists $\alpha \in \mathbb{R}$ such that
	\[
	\alpha\ge 0, \alpha^2=2
	.\] 
\end{theorem}
\begin{proof}
Let $S= \{s\ge 0:s^2<2\} $ Since  $S$ is nonempty and bounded from above by $2$, by Supremum Property, there is a supremum $\alpha = \operatorname{sup} S$. We claim that $\alpha^2 = 2$.

Suppose not, then by Trichotomy, either $\alpha^2 >2 $ or $\alpha^2 <2$. Suppose $\alpha^2 <2$. Then we can find large $n$ such that  $\alpha + \frac{1}{n} < 2$. Suppose $\alpha^2 > 2$, we can find $\left( \alpha - \frac{1}{m} \right)  > 2$.
In both cases $\alpha \neq \operatorname{sup} S$.
\end{proof}

\begin{remark}
	We used $\frac{1}{n}$ instead of $\varepsilon \in \mathbb{R}$ in the proof. This is not strictly necessary. However, if we construct the reals using Dedekind cuts, we have to work with rationals.
\end{remark}

\subsection{Cantor Set and Nested Intervals}

\begin{definition}[interval]
	For $a,b$ in $\mathbb{R}$, define 
	\begin{itemize}
		\item (closed interval) $[a,b] = \{x \in \mathbb{R}: a \le x\le b\} $ 
		\item (open interval)$(a,b) = \{x \in \mathbb{R}: a < x <  b$\} 
		\item  (semi-open / semi-closed) $(a,b] = \{x \in \mathbb{R}\} : a < x \le  b$
		\item  $[a,b) = \{x \in \mathbb{R}\} : a \le  x  < b$
		\item (semi-infinite / rays) $[a,\infty ) = \{x \in \mathbb{R}: x \ge a\} $
		\item  $(a,\infty ) = \{x \in \mathbb{R}: x > a\} $ 
		\item (infinite) $(-\infty , \infty) = \mathbb{R}$
	\end{itemize}
\end{definition}

\begin{theorem}[Nested Interval Property]
	Suppose we have a sequence of closed finite nonempty intervals $I_1, I_2, \ldots$, such that \[
	I_1 \supset I_2 \supset \ldots \supset I_n \supset \ldots
\]
Then there exists $x \in \mathbb{R}$ such that $x \in I_n$ for any $n \in  \mathbb{N}$. In other words, $\bigcap_{n \in \mathbb{N}} I_n \neq  \emptyset  $
\end{theorem}
\begin{proof}
	Let $I_n = \left[ a_n, b_n \right] $. Suppose WOLG that $a_n \le b_n$ then we have
\begin{align*}
	a_1\le a_2\le \ldots\le a_n \le \ldots \\
	b_1\ge b_2\ge \ldots\ge b_n \ge \ldots 
.\end{align*}
\textbf{Claim.} $a_m \le  b_k $ for any $m, k \in \mathbb{N}$. \textit{Proof:} Let $n = \max \{m,k\} $. Then $a_m \le a_n \le b_n \le b_k$.

$b_k$ is an upper bound for $A = \{a_n : n \in \mathbb{N}\} $. $A\neq \emptyset $, so there exists, by supremum property, some $\xi = \operatorname{sup} A$. By definition of supremum, $a_n \le  \xi$ for any $ n \in \mathbb{N}$.

On the other hand, $b_k$ is an upper bound, so $\xi \le b_k$ for any $k \in \mathbb{N}$. Now $a_n \le \xi \le b_n$, so $\xi \in I_n$.
\end{proof}

